\documentclass[11pt,a4paper]{article}
\usepackage[utf8]{inputenc}
\usepackage[T1]{fontenc}
\usepackage{lmodern}
\usepackage{amsmath}
\usepackage[margin=2cm]{geometry}
\usepackage{enumitem}
\pagestyle{plain}
\setlength{\parindent}{0pt}
\setlength{\parskip}{0pt}
\begin{document}
{\Large\bfseries Tarea de sales ácidas y básicas}
\vspace{12pt}
\begin{enumerate}[label=\textbf{\arabic*.},leftmargin=*,itemsep=11pt,topsep=0pt]
\item \begin{minipage}[t]{\linewidth}
¿Cuál es la fórmula del hidrogenocarbonato de calcio?

\vspace{5pt}
\begin{tabular}{@{}p{0.47\linewidth}p{0.47\linewidth}@{}}
a)~$\mathrm{Ca_{2}HCO_{3}}$ & b)~$\mathrm{CaHCO_{3}}$ \\
c)~$\mathrm{Ca(HCO_{3})_{2}}$ & d)~$\mathrm{CaH_{2}CO_{3}}$ \\
\end{tabular}
\end{minipage}

\item \begin{minipage}[t]{\linewidth}
¿Cuál es el nombre tradicional del $\mathrm{Na_{2}HPO_{4}}$?

\vspace{5pt}
\begin{tabular}{@{}p{0.47\linewidth}p{0.47\linewidth}@{}}
a)~fosfato básico de sodio & b)~fosfato ácido de sodio \\
c)~fosfato diácido de sodio & d)~fosfito ácido de sodio \\
\end{tabular}
\end{minipage}

\item \begin{minipage}[t]{\linewidth}
¿Cuál es el nombre sistemático del $\mathrm{KHCO_{3}}$?

\vspace{5pt}
\begin{tabular}{@{}p{0.47\linewidth}p{0.47\linewidth}@{}}
a)~hidrogenodioxocarbonato (II) de potasio & b)~trioxocarbonato (IV) de potasio \\
c)~hidrogenotetraoxocarbonato (IV) de potasio & d)~hidrogenotrioxocarbonato (IV) de potasio \\
\end{tabular}
\end{minipage}

\item \begin{minipage}[t]{\linewidth}
¿Cómo se forma una sal básica?

\vspace{5pt}
\begin{tabular}{@{}p{0.47\linewidth}p{0.47\linewidth}@{}}
a)~un óxido se une con agua & b)~un ácido neutraliza solo parte de los OH \\
c)~un metal reemplaza parte de los H de un ácido & d)~un no metal se une con hidrógeno \\
\end{tabular}
\end{minipage}

\item \begin{minipage}[t]{\linewidth}
¿Cuál es la fórmula de la sal básica de plomo (II) con un OH y un nitrato?

\vspace{5pt}
\begin{tabular}{@{}p{0.47\linewidth}p{0.47\linewidth}@{}}
a)~$\mathrm{PbOH(NO_{3})_{2}}$ & b)~$\mathrm{Pb(OH)_{2}NO_{3}}$ \\
c)~$\mathrm{Pb(OH)NO_{3}}$ & d)~$\mathrm{Pb_{2}(OH)NO_{3}}$ \\
\end{tabular}
\end{minipage}

\item \begin{minipage}[t]{\linewidth}
¿Cuál es el nombre sistemático del $\mathrm{Fe(OH)CO_{3}}$?

\vspace{5pt}
\begin{tabular}{@{}p{0.47\linewidth}p{0.47\linewidth}@{}}
a)~hidroxitrioxocarbonato (IV) de hierro (III) & b)~carbonato básico de hierro (II) \\
c)~hidroxidioxocarbonato (IV) de hierro (III) & d)~hidroxitrioxocarbonato (IV) de hierro (II) \\
\end{tabular}
\end{minipage}

\item \begin{minipage}[t]{\linewidth}
¿Cuál es el nombre tradicional del $\mathrm{Al(OH)_{2}NO_{3}}$?

\vspace{5pt}
\begin{tabular}{@{}p{0.47\linewidth}p{0.47\linewidth}@{}}
a)~nitrato dibásico de aluminio & b)~nitrato básico de aluminio \\
c)~nitrato diácido de aluminio & d)~nitrito dibásico de aluminio \\
\end{tabular}
\end{minipage}

\item \begin{minipage}[t]{\linewidth}
¿Qué tipo de compuesto es el $\mathrm{KClO_{3}}$?

\vspace{5pt}
\begin{tabular}{@{}p{0.47\linewidth}p{0.47\linewidth}@{}}
a)~hidróxido & b)~sal ácida \\
c)~oxiácido & d)~oxisal neutra \\
\end{tabular}
\end{minipage}

\item \begin{minipage}[t]{\linewidth}
¿Cuál es el nombre stock del $\mathrm{Ni_{2}(OH)_{4}SO_{3}}$?

\vspace{5pt}
\begin{tabular}{@{}p{0.47\linewidth}p{0.47\linewidth}@{}}
a)~tetrahidróxido-sulfito de níquel (III) & b)~dihidróxido-sulfito de níquel (III) \\
c)~tetrahidróxido-sulfato de níquel (III) & d)~tetrahidróxido-sulfito de níquel (II) \\
\end{tabular}
\end{minipage}

\item \begin{minipage}[t]{\linewidth}
Por su fórmula, ¿qué tipo de compuesto es el tetraoxosulfato (VI) de hidrógeno?

\vspace{5pt}
\begin{tabular}{@{}p{0.47\linewidth}p{0.47\linewidth}@{}}
a)~oxisal neutra & b)~oxiácido \\
c)~sal binaria & d)~sal ácida \\
\end{tabular}
\end{minipage}

\end{enumerate}
\end{document}
